% Part: normal-modal-logic
% Chapter: axioms-systems
% Section: logics-proofs

\documentclass[../../../include/open-logic-section]{subfiles}

\begin{document}

\olfileid[fr]{nml}{prf}{prf}

\olsection{\usetoken{P}{derivation} et systèmes modaux}

Commençons par définir ce qu'est une !!{derivation} pour les logiques
modales normales. Schématiquement, il s'agit d'une suite de
!!{formula}s dont chaque !!{element} est soit une instance par
substitution de l'un des \emph{axiomes}, soit une conséquence des
!!{element}s précédents par quelques règles d'inférence. Dans les
logiques modales normales, toutes les instances des
tautologies\iftag{prvDiamond}{, de \Ax{K} et de~\Dual{}}{ et de~\Ax{K}}
comptent comme axiomes. On obtient ainsi le système modal~$\Log{K}$,
la plus petite logique modale normale. On peut ajouter d'autres
axiomes pour obtenir d'autres systèmes. Les règles restent toujours
le modus ponens~\MP{} et la nécessitation~\Nec.

\begin{defn}
  Étant donné un système modal $\Log{K} !A_1 \dots !A_n$ et une
  !!{formula}~$!B$, nous disons que $!B$ est
  \emph{!!{derivable}} dans $\Log{K} !A_1 \dots !A_n$, ce que nous
  notons $\Log{K} !A_1 \dots !A_n \Proves !B$, si et seulement s'il
  existe des !!{formula}s $!C_1$, \dots, $!C_k$ telles que
  $!C_k = !B$ et que chaque $!C_i$ est une instance tautologique,
  une instance de l'un des axiomes $\Ax{K}$,\iftag{prvDiamond}{
  $\Dual$,}{} $!A_1$, \dots, $!A_n$, ou découle de !!{formula}s
  précédentes par les règles \MP{} ou~\Nec.
\end{defn}

La proposition suivante permet de montrer que $!B \in \Sigma$ en
exhibant une !!{derivation} de~$!B$ dans~$\Sigma$.

\begin{prop}
  $\Log{K} !A_1 \dots !A_n = \Setabs{!B}{\Log{K} !A_1 \dots !A_n \Proves !B}$.
\end{prop}

\begin{proof}
  Par récurrence sur la longueur des !!{derivation}s, montrons que
  $\Setabs{!B}{\Log{K} !A_1 \dots !A_n \Proves !B} \subseteq
  \Log{K} !A_1 \dots !A_n$.

  Si la !!{derivation} de~$!B$ a la longueur~$1$, elle ne contient
  qu'une !!{formula}. Celle-ci ne peut découler de formules
  précédentes par \MP{} ou \Nec; elle est donc une instance
  tautologique, une instance de \Ax{K},\iftag{prvDiamond}{ de~\Dual{},}{}
  ou une instance de l'un de $!A_1$, \dots,~$!A_n$. Or
  $\Log{K}!A_1\dots!A_n$ contient aussi ces instances; ainsi
  $!B \in \Log{K}!A_1 \dots !A_n$.

  Si la !!{derivation} de $!B$ a une longueur~$> 1$, et si sa
  dernière ligne est un axiome, le raisonnement précédent s'applique.
  Sinon, cette ligne découle par \MP{} ou~\Nec{} de lignes
  antérieures. Si $!B$ découle de $!C$ et de
  $!C \lif !B$ par \MP, l'hypothèse de récurrence donne
  $!C$ et $!C \lif !B \in \Log{K}!A_1 \dots !A_n$. Toute logique
  modale étant close par modus ponens,
  $!B \in \Log{K}!A_1 \dots !A_n$. Si $!B \equiv \Box !C$ découle de
  $!C$ par \Nec, l'hypothèse de récurrence donne
  $!C \in \Log{K}!A_1 \dots !A_n$. Toute logique modale normale
  étant close par \Nec, nous avons aussi
  $!B \in \Log{K}!A_1\dots!A_n$.

  Pour l'inclusion réciproque, montrons que
  $\Sigma = \Setabs{!B}{\Log{K} !A_1 \dots !A_n \Proves !B }$ est une
  logique modale normale contenant toutes les instances de $!A_1$,
  \dots, $!A_n$. Chacune de ces instances a, par définition, une
  !!{derivation} d'une seule ligne. Le système $\Log{K} !A_1
  \dots !A_n$ est la plus petite de ces logiques.
  \begin{enumerate}
    \item Toute tautologie~$!B$ est une instance tautologique; donc
      $\Log{K}!A_1\dots!A_n \Proves !B$ et $\Sigma$ contient toutes
      les tautologies.
    \item Si $\Log{K}!A_1\dots!A_n \Proves !C$ et
      $\Log{K}!A_1\dots!A_n \Proves !C \lif !B$, alors
      $\Log{K}!A_1\dots!A_n \Proves !B$ : il suffit de réunir la
      !!{derivation} de $!C$ et celle de $!C \lif !B$, puis d'ajouter
      la ligne~$!B$, justifiée par \MP{}. Ainsi, $\Sigma$ est clos
      par modus ponens.
    \item Si $!B$ possède une !!{derivation}, chacune de ses instances
      par substitution en possède aussi une : appliquons la
      substitution à chaque !!{formula} de la !!{derivation}.
      (Exercice : démontrer par récurrence sur la longueur des
      !!{derivation}s que le résultat est encore une !!{derivation}
      correcte.) Ainsi, $\Sigma$ est clos par substitution uniforme.
      (Nous avons établi toutes les conditions pour que $\Sigma$
      soit une logique modale.)
    \item Toute instance de \Ax{K} est un axiome, donc elle appartient
      à~$\Sigma$.
    \tagitem{prvDiamond}{Toute instance de \Dual{} est également un
      axiome, donc elle appartient à~$\Sigma$.}{}
    \item Si $\Log{K}!A_1\dots!A_n \Proves !C$, la ligne
      supplémentaire~$\Box !C$ est justifiée par~\Nec. Par
      conséquent, $\Sigma$ est clos par~\Nec, donc normal.
  \end{enumerate}
\end{proof}

\end{document}
